% Original problem figures and added solution diagrams for Section 3.3.
\tikzset{
  solbody/.style={draw=Ink,line width=0.8pt},
  solforce/.style={-{Latex[length=2.0mm]},draw=DeepBlue,line width=0.9pt},
  solgeom/.style={draw=TextGray,densely dashed,line width=0.55pt},
  solflow/.style={-{Latex[length=1.8mm]},draw=Blue,line width=0.75pt},
  sollabel/.style={font=\scriptsize,fill=none,inner sep=1pt},
  solpanel/.style={font=\scriptsize\bfseries,text=TextGray}
}
\newcommand{\diagcaption}[2]{%
 \refstepcounter{figure}\par\nopagebreak[4]\smallskip
 {\centering\scriptsize\color{TextGray}图~\thefigure\quad #2\par}\smallskip}
% Shared geometry for repeated momentum problems. Captions remain at call sites.
\newcommand{\momentumRain}{%
\begin{tikzpicture}[scale=.78]
  \fill[Blue!8] (-1.2,-.65) rectangle (1.2,-.12);
  \draw[solbody] (-1.2,.7)--(-1.2,-.65)--(1.2,-.65)--(1.2,.7);
  \foreach \x in {-.9,-.45,0,.45,.9}{\draw[solflow] (\x,1.55)--(\x,.86);}
  \node[sollabel] at (0,2.05) {$q=\rho S\Delta h/\tau$，雨水水平速度为零};
  \node[sollabel] at (0,-1.15) {$m(t)=M+qt$};
  \draw[solforce] (1.2,-.38)--(2.55,-.38) node[midway,below=3pt,sollabel] {$F$};
  \draw[solflow] (1.6,.52)--(2.8,.52) node[midway,above=3pt,sollabel] {$v$};
\end{tikzpicture}}

\newcommand{\momentumPickup}{%
\begin{tikzpicture}[scale=.8]
  \draw[solbody] (-2,0)--(2.5,0);
  \draw[DeepBlue,line width=1.2pt] (-1.8,0)--(-.15,0)
    .. controls (-.05,0) and (0,.05) .. (0,.15)--(0,2.4);
  \fill[Ink] (0,2.4) circle (2pt);
  \draw[solforce] (0,2.4)--(0,3.3) node[left=3pt,sollabel] {$F$};
  \draw[solforce] (-.65,1.65)--(-.65,.7) node[midway,left=3pt,sollabel] {$\lambda xg$};
  \draw[solflow] (.65,1.1)--(.65,2.1) node[midway,right=3pt,sollabel] {$v$};
  \draw[solgeom] (0,2.4)--(1.8,2.4);
  \draw[<->,draw=TextGray] (1.8,0)--(1.8,2.4) node[midway,right=3pt,sollabel] {$x$};
  \node[sollabel] at (-1.0,-.38) {静止绳段};
  \node[sollabel] at (0,-.95) {$\dot m_{\rm mov}=\lambda v$，新增动量率 $\lambda v^2$};
\end{tikzpicture}}

\newcommand{\momentumFallingRope}{%
\begin{tikzpicture}[scale=.78]
  \draw[solbody] (-.45,3)--(.45,3);
  \foreach \x in {-.4,-.2,0,.2,.4}{\draw[TextGray] (\x,3)--++(.12,.14);}
  \draw[solgeom] (0,3)--(0,2.1);
  \draw[DeepBlue,line width=1.1pt] (0,2.1)--(0,.28)
    .. controls (0,.08) and (-.15,.08) .. (-.3,.08);
  \draw[DeepBlue,line width=1.1pt,domain=-.3:1.6,samples=120]
    plot ({\x},{.08+.08*sin((\x+.3)*720)});
  \draw[solbody] (-.75,0)--(1.9,0);
  \draw[solgeom] (-.7,3)--(0,3) (-.7,2.1)--(0,2.1);
  \draw[<->,draw=TextGray] (-.7,2.1)--(-.7,3) node[midway,left=3pt,sollabel] {$x$};
  \draw[solflow] (.55,1.8)--(.55,.9) node[midway,right=3pt,sollabel] {$v$};
  \node[sollabel,anchor=west] at (.65,2.8) {虚线：上端已下降的距离};
  \node[sollabel,anchor=west] at (1.15,1.5) {$l-x$：仍在下落的绳长};
  \node[sollabel] at (.4,-.4) {秤盘；堆积绳段长度为 $x$};
\end{tikzpicture}}

\newcommand{\momentumRopeLoad}{%
\begin{tikzpicture}[scale=.78]
  \draw[DeepBlue,line width=1.1pt] (0,2)--(0,.28)
    .. controls (0,.08) and (-.15,.08) .. (-.3,.08);
  \draw[DeepBlue,line width=1.1pt,domain=-.3:1.8,samples=120]
    plot ({\x},{.08+.08*sin((\x+.3)*720)});
  \draw[solbody] (-.65,0)--(2.05,0);
  \draw[solflow] (.55,1.7)--(.55,.85) node[midway,right=3pt,sollabel] {$v$};
  \draw[solforce] (.45,0)--(.45,-1.0) node[below=3pt,sollabel] {$\lambda xg$};
  \draw[solforce] (1.6,0)--(1.6,-1.0) node[below=3pt,sollabel] {$\lambda v^2$};
  \node[sollabel,anchor=west,align=left] at (2.7,1.15)
    {秤盘承受的向下作用\\$N=\lambda xg+\lambda v^2$\\第二项来自使绳元停止};
\end{tikzpicture}}

\newcommand{\momentumAccretion}{%
\begin{tikzpicture}[scale=.8]
  \draw[solbody,fill=SoftGray] (-.8,0) rectangle (.8,.9);
  \draw[solforce] (-.2,.45)--(-.2,-.85) node[below=3pt,sollabel] {$mg$};
  \draw[solflow] (1.25,.7)--(1.25,-.25) node[midway,right=3pt,sollabel] {$v$};
  \foreach \x in {-1.3,-.85,-.4,.05,.5,.95,1.4}{\fill[TextGray] (\x,-1.55) circle (1.5pt);}
  \node[sollabel] at (0,-1.95) {地面系中静止的介质};
  \draw[<->,TextGray] (-.8,1.3)--(.8,1.3) node[midway,above=3pt,sollabel] {截面积 $S$};
  \node[sollabel,anchor=west,align=left] at (2.3,.8)
    {$\dot m=\rho Sv$\\$\dd(mv)/\dd t=mg$\\$m\dot v=mg-\rho Sv^2$};
\end{tikzpicture}}

\newcommand{\momentumDroplet}{%
\begin{tikzpicture}[scale=.8]
  \draw[solbody,fill=Blue!8] (0,.5) circle (.7);
  \fill[Ink] (0,.5) circle (1.5pt);
  \draw[solgeom] (0,.5)--(.7,.5) node[midway,above=4pt,sollabel] {$r$};
  \draw[solforce] (0,.5)--(0,-1.0) node[below=3pt,sollabel] {$mg$};
  \draw[solflow] (1.25,.75)--(1.25,-.25) node[midway,right=3pt,sollabel] {$v$};
  \foreach \x in {-1.3,-.85,-.4,.05,.5,.95,1.4}{\fill[TextGray] (\x,-1.65) circle (1.5pt);}
  \node[sollabel] at (0,-2.05) {地面系中静止的雾滴};
  \node[sollabel,anchor=west,align=left] at (2.3,.8)
    {$S=\pi r^2$\\$\dot m=\rho_f\pi r^2v$\\$\dot r=\rho_f v/(4\rho_w)$};
\end{tikzpicture}}


\newcommand{\soldiagThreeTwentyFour}{%
\begin{center}\begin{minipage}{.97\linewidth}\centering
\momentumRain
\diagcaption{解图}{雨水只增加质量，不带入水平动量；外力使系统水平动量按 $Ft$ 增长。}\label{ch07:fig:soldiagThreeTwentyFour}\end{minipage}\end{center}}

\newcommand{\soldiagThreeTwentyFive}{%
\begin{center}\begin{minipage}{.97\linewidth}\centering
\momentumPickup
\diagcaption{解图}{提力既承担运动绳段的重量，也为新进入运动段的绳元提供动量。}\label{ch07:fig:soldiagThreeTwentyFive}\end{minipage}\end{center}}

\newcommand{\probdiagThreeTwentySix}{%
\begin{center}\begin{minipage}{.97\linewidth}\centering
\momentumFallingRope
\diagcaption{题图}{柔绳上端释放后下落，绳元依次落到秤盘上并停止。}\label{ch07:fig:probdiagThreeTwentySix}\end{minipage}\end{center}}

\newcommand{\soldiagThreeTwentySix}{%
\begin{center}\begin{minipage}{.97\linewidth}\centering
\momentumRopeLoad
\diagcaption{解图}{秤盘既承受已落绳段的重量，也连续吸收下落绳元的动量。}\label{ch07:fig:soldiagThreeTwentySix}\end{minipage}\end{center}}

\newcommand{\soldiagThreeTwentySeven}{%
\begin{center}\begin{minipage}{.97\linewidth}\centering
\momentumAccretion
\diagcaption{解图}{静止介质被吸收后要加速到 $v$，故出现随 $v^2$ 增长的动量流项。}\label{ch07:fig:soldiagThreeTwentySeven}\end{minipage}\end{center}}

\newcommand{\soldiagThreeTwentyEight}{%
\begin{center}\begin{minipage}{.97\linewidth}\centering
\momentumDroplet
\diagcaption{解图}{半径、迎风面积和质量共同增长，速度必须与它们联立求解。}\label{ch07:fig:soldiagThreeTwentyEight}\end{minipage}\end{center}}

\newcommand{\rocketMicroDiagram}{%
\begin{center}\begin{minipage}{.97\linewidth}\centering
\begin{tikzpicture}[scale=.77]
 \node[solpanel] at (-2.8,1.5) {喷气前}; \node[solpanel] at (2.6,1.5) {喷气后};
 \draw[solbody,fill=SoftGray] (-4,-.35) rectangle (-2,.35);
 \node[sollabel] at (-3,0) {$M$};
 \draw[solflow] (-2,0)--(-.7,0) node[midway,above=4pt,sollabel] {$v$};
 \draw[solbody,fill=SoftGray] (1.4,-.35) rectangle (3.4,.35);
 \node[sollabel] at (2.4,0) {$M+\Delta M$};
 \draw[solflow] (3.4,0)--(4.7,0) node[midway,above=4pt,sollabel] {$v+\Delta v$};
 \foreach \x in {-.2,.2,.6}{\fill[TextGray] (\x,0) circle (2pt);}
 \node[sollabel] at (.2,.75) {燃气 $-\Delta M$};
 \draw[solflow] (.65,-.95)--(-.65,-.95) node[midway,below=4pt,sollabel] {$-u$：相对火箭};
 \node[sollabel] at (.2,-1.8) {燃气对地速度为 $v-u$，可以向左、为零或向右};
\end{tikzpicture}
\diagcaption{示意图}{在地面系中，喷出燃气的速度是 $v-u$；$\Delta M<0$。}\label{ch07:fig:rocketMicroDiagram}\end{minipage}\end{center}}

\newcommand{\probdiagThreeTwentyNine}{%
\begin{center}\begin{minipage}{.97\linewidth}\centering
\begin{tikzpicture}[xscale=.82,yscale=.70]
  \draw[axis] (0,0)--(5.6,0) node[right,sollabel] {火箭质量};
  \draw[axis] (0,0)--(0,3.8) node[above,sollabel] {火箭速度};
  \draw[solgeom] (4.45,0)--(4.45,1.15)--(0,1.15);
  \fill[Ink] (4.45,1.15) circle (2.2pt);
  \node[sollabel,below] at (4.45,0) {$M$};
  \node[sollabel,left] at (0,1.15) {$v$};
  \node[sollabel,above right] at (4.45,1.15) {当前状态};
\end{tikzpicture}
\diagcaption{题图}{曲线应从当前状态 $(M,v)$ 出发，并标明燃烧过程的方向。}\label{ch07:fig:probdiagThreeTwentyNine}\end{minipage}\end{center}}

\newcommand{\soldiagThreeTwentyNine}{%
\begin{center}\begin{minipage}{.97\linewidth}\centering
\begin{tikzpicture}[scale=.82]
 \draw[axis] (0,0)--(5.2,0) node[right,sollabel] {$\mathcal M$};
 \draw[axis] (0,0)--(0,3.7) node[above,sollabel] {$V$};
 \draw[curve,domain=.6:4.4,samples=80] plot ({\x},{1+ln(4.4/\x)});
 \draw[curve,-{Latex[length=2mm]},domain=3.4:2.4,samples=20] plot ({\x},{1+ln(4.4/\x)});
 \draw[solgeom] (0,1)--(4.4,1)--(4.4,0);
 \node[sollabel,left=4pt] at (0,1) {$v$}; \node[sollabel,below=4pt] at (4.4,0) {$M$};
 \node[sollabel,anchor=west] at (2.5,3.2) {$V=v+u\ln(M/\mathcal M)$};
 \node[sollabel,anchor=west] at (2.5,2.6) {燃烧方向：沿曲线向左上方};
\end{tikzpicture}
\diagcaption{解图}{横轴质量减小时，火箭沿对数曲线向左上方运动。}\label{ch07:fig:soldiagThreeTwentyNine}\end{minipage}\end{center}}

\newcommand{\soldiagThreeThirty}{%
\begin{center}\begin{minipage}{.97\linewidth}\centering
\begin{tikzpicture}[scale=.82]
 \draw[axis] (0,0)--(5.8,0) node[right,sollabel] {$\mathcal M$};
 \draw[axis] (0,0)--(0,3.8) node[above,sollabel] {$u/\mathcal M$};
 \path[fill=Blue!12] (1.5,0)--plot[domain=1.5:4.6,samples=50] (\x,{2.1/\x})--(4.6,0)--cycle;
 \draw[curve,domain=.6:5.1,samples=80] plot (\x,{2.1/\x});
 \draw[solgeom] (1.5,0)--(1.5,1.4) (4.6,0)--(4.6,{2.1/4.6});
 \node[sollabel,below=4pt] at (1.5,0) {$M$}; \node[sollabel,below=4pt] at (4.6,0) {$M_0$};
 \node[sollabel] (L) at (3.8,2.5) {阴影面积 $u\ln(M_0/M)$};
 \draw[TextGray] (L.south)--(3,.4);
 \draw[solflow] (4.6,-.65)--(1.5,-.65) node[midway,below=4pt,sollabel] {质量减少的方向};
\end{tikzpicture}
\diagcaption{解图}{对 $\dd v=-u\,\dd M/M$ 积分，速度增量等于 $u/M$ 曲线下对应的面积。}\label{ch07:fig:soldiagThreeThirty}\end{minipage}\end{center}}
